% GENERATED FILE. Edit canonical lessons, then run npm run generate:books.
% canonical-source-hash: fnv1a64:5b2329271bb06eee
% canonical-lessons: GE-C276-umfrage, GE-C276-statistik, GE-C276-anteil, GE-C276-zunehmen, GE-C276-sinken

\chapter{Zahlen in den Nachrichten}
\label{ch:ge-zahlen-in-den-nachrichten}
\hlchaptermodality{276}

\begin{chapteropening}
\textbf{By the end of this chapter:} \emph{I can report numbers in the news: a survey, statistics, a share, to increase and to fall.}

Complete the last lesson of chapter 276: I can report numbers in the news: a survey, statistics, a share, to increase and to fall.
\end{chapteropening}

\section[die Umfrage]{die Umfrage — the survey, the poll}

\label{lesson:GE-C276-umfrage}



\begin{quote}
Before the new one: say the German for to suspect, then the German for presumably.
\end{quote}

\subsection*{You'll want to know: die Umfrage}
\textbf{die Umfrage} — \textquotedblleft{}the survey, the poll\textquotedblright{}.

\begin{quote}
Bei einer Umfrage sagten viele Leute Nein.
\end{quote}

\begin{quote}
\emph{In a survey many people said no.}
\end{quote}

Plural: \emph{die Umfragen}.

\begin{cousinweb}
\emph{Umfrage} joins \emph{um} \textquotedblleft{}around\textquotedblright{} and \emph{Frage} \textquotedblleft{}question\textquotedblright{}, both inherited Germanic words: a question asked all around. EtymDB lists no parent for the compound.
\end{cousinweb}

\begin{grammarlens}[title={What you've built}]
One more word for telling, explaining and describing at B1.
\end{grammarlens}

\subsection*{Your turn}
\begin{itemize}
\raggedright
  \item \emph{Say it:} \emph{die Umfrage}
  \item \emph{Say it:} \emph{die Umfrage}, once more
  \item \emph{Say it:} a sentence of your own with \emph{die Umfrage}
  \item \emph{Recall:} say the German for to suspect, then the German for presumably, then say \emph{die Umfrage} again
\end{itemize}

\begin{tcolorbox}[breakable,colback=teal!4,colframe=teal!35!black,title={Before you move on}]
What does die Umfrage mean? (\textquotedblleft{}The survey, the poll\textquotedblright{}.) Say it once more.
\end{tcolorbox}
\section[die Statistik]{die Statistik — the statistics; a statistic}

\label{lesson:GE-C276-statistik}



\begin{quote}
Before the new one: say the German for presumably, then the German for the survey.
\end{quote}

\subsection*{You'll want to know: die Statistik}
\textbf{die Statistik} — \textquotedblleft{}the statistics; a statistic\textquotedblright{}.

\begin{quote}
Die Statistik zeigt mehr Unfälle im Winter.
\end{quote}

\begin{quote}
\emph{The statistics show more accidents in winter.}
\end{quote}

Plural: \emph{die Statistiken}. Stress on the middle: \emph{Sta-TIS-tik}.

\begin{cousinweb}
\emph{Statistik} is usually traced to Neo-Latin \emph{statisticus} \textquotedblleft{}of the state\textquotedblright{}, built on Latin \emph{status} \textquotedblleft{}state, condition\textquotedblright{} (Lewis \& Short): first, facts about a state. EtymDB lists no parent.
\end{cousinweb}

\begin{grammarlens}[title={What you've built}]
One more word for telling, explaining and describing at B1.
\end{grammarlens}

\subsection*{Your turn}
\begin{itemize}
\raggedright
  \item \emph{Say it:} \emph{die Statistik}
  \item \emph{Say it:} \emph{die Statistik}, once more
  \item \emph{Say it:} a sentence of your own with \emph{die Statistik}
  \item \emph{Recall:} say the German for presumably, then the German for the survey, then say \emph{die Statistik} again
\end{itemize}

\begin{tcolorbox}[breakable,colback=teal!4,colframe=teal!35!black,title={Before you move on}]
What does die Statistik mean? (\textquotedblleft{}The statistics; a statistic\textquotedblright{}.) Say it once more.
\end{tcolorbox}
\section[der Anteil]{der Anteil — the share, the proportion}

\label{lesson:GE-C276-anteil}



\begin{quote}
Before the new one: say the German for the survey, then the German for the statistics.
\end{quote}

\subsection*{You'll want to know: der Anteil}
\textbf{der Anteil} — \textquotedblleft{}the share, the proportion\textquotedblright{}.

\begin{quote}
Der Anteil der Kinder war sehr hoch.
\end{quote}

\begin{quote}
\emph{The share of children was very high.}
\end{quote}

Plural: \emph{die Anteile}. \emph{Anteil nehmen} — to sympathise.

\begin{cousinweb}
\emph{Anteil} is built on \emph{Teil} \textquotedblleft{}part\textquotedblright{}, from Old High German \emph{teil}, Proto-Germanic \emph{dailą} \textquotedblleft{}part, deal\textquotedblright{} (EtymDB), with \emph{an}.
\end{cousinweb}

\begin{grammarlens}[title={What you've built}]
One more word for telling, explaining and describing at B1.
\end{grammarlens}

\subsection*{Your turn}
\begin{itemize}
\raggedright
  \item \emph{Say it:} \emph{der Anteil}
  \item \emph{Say it:} \emph{der Anteil}, once more
  \item \emph{Say it:} a sentence of your own with \emph{der Anteil}
  \item \emph{Recall:} say the German for the survey, then the German for the statistics, then say \emph{der Anteil} again
\end{itemize}

\begin{tcolorbox}[breakable,colback=teal!4,colframe=teal!35!black,title={Before you move on}]
What does der Anteil mean? (\textquotedblleft{}The share, the proportion\textquotedblright{}.) Say it once more.
\end{tcolorbox}
\section[zunehmen]{zunehmen — to increase; to put on weight}

\label{lesson:GE-C276-zunehmen}



\begin{quote}
Before the new one: say the German for the statistics, then the German for the share.
\end{quote}

\subsection*{You'll want to know: zunehmen}
\textbf{zunehmen} — \textquotedblleft{}to increase; to put on weight\textquotedblright{}.

\begin{quote}
Der Verkehr nahm am Abend zu.
\end{quote}

\begin{quote}
\emph{The traffic increased in the evening.}
\end{quote}

Separable: \emph{es nimmt zu}. For a person, to put on weight: \emph{Er hat zugenommen}.

\begin{cousinweb}
\emph{zunehmen} is built on \emph{nehmen} \textquotedblleft{}to take\textquotedblright{}, from Old High German \emph{neman}, Proto-Germanic \emph{nemaną} (EtymDB), with \emph{zu}: to take on more.
\end{cousinweb}

\begin{grammarlens}[title={What you've built}]
One more word for telling, explaining and describing at B1.
\end{grammarlens}

\subsection*{Your turn}
\begin{itemize}
\raggedright
  \item \emph{Say it:} \emph{zunehmen}
  \item \emph{Say it:} \emph{zunehmen}, once more
  \item \emph{Say it:} a sentence of your own with \emph{zunehmen}
  \item \emph{Recall:} say the German for the statistics, then the German for the share, then say \emph{zunehmen} again
\end{itemize}

\begin{tcolorbox}[breakable,colback=teal!4,colframe=teal!35!black,title={Before you move on}]
What does zunehmen mean? (\textquotedblleft{}To increase; to put on weight\textquotedblright{}.) Say it once more.
\end{tcolorbox}
\section[sinken]{sinken — to sink, to fall}

\label{lesson:GE-C276-sinken}



\begin{quote}
Before the new one: say the German for the share, then the German for to increase.
\end{quote}

\subsection*{You'll want to know: sinken}
\textbf{sinken} — \textquotedblleft{}to sink, to fall\textquotedblright{}.

\begin{quote}
Die Preise sanken im Sommer.
\end{quote}

\begin{quote}
\emph{Prices fell in the summer.}
\end{quote}

Strong verb with \emph{sein}: \emph{sank}, \emph{ist gesunken}.

\begin{cousinweb}
\emph{sinken} is an inherited Germanic verb, from Old High German \emph{sinkan} (EtymDB), and EtymDB lists Low German \emph{sinken} and English \emph{sink} among its cousins.
\end{cousinweb}

\begin{grammarlens}[title={What you've built}]
One more word for telling, explaining and describing at B1.
\end{grammarlens}

\subsection*{Your turn}
\begin{itemize}
\raggedright
  \item \emph{Say it:} \emph{sinken}
  \item \emph{Say it:} \emph{sinken}, once more
  \item \emph{Say it:} a sentence of your own with \emph{sinken}
  \item \emph{Recall:} say the German for the share, then the German for to increase, then say \emph{sinken} again
\end{itemize}

\begin{tcolorbox}[breakable,colback=teal!4,colframe=teal!35!black,title={Before you move on}]
What does sinken mean? (\textquotedblleft{}To sink, to fall\textquotedblright{}.) Say it once more.
\end{tcolorbox}
